% File di appoggio: sezioni "Santo" e "Alleluia e Lode" tolte dal libretto (10/11/2026).
% Per rimetterle: % File di appoggio: sezioni "Santo" e "Alleluia e Lode" tolte dal libretto (10/11/2026).
% Per rimetterle: % File di appoggio: sezioni "Santo" e "Alleluia e Lode" tolte dal libretto (10/11/2026).
% Per rimetterle: % File di appoggio: sezioni "Santo" e "Alleluia e Lode" tolte dal libretto (10/11/2026).
% Per rimetterle: \input{appoggio/santi_alleluia} in libretto.tex (o ricopiare le due parti in canti/).

% ===================== SANTO =====================
% !TEX root = ../libretto.tex
% Sezione: Santo
% Nei brani a due cori (Uomini/Donne) le parti sono state unite in un solo testo.
\momento{Santo}
\begin{songs}{titoli}

% Fonte: LibrettoAccordi.pdf (Canti Coro), H10.
% Nell'elenco: Santo, Ordinario.
\beginsong{Osanna eh}[by={RnS}]
\beginchorus
Santo! Santo! Osanna! \rep{2}
Osanna eh! Osanna eh! Osanna a Cristo Signor! \rep{2}
\endchorus
\beginverse*
I cieli e la terra, o Signore, sono pieni di te. \rep{2}
\endverse
\beginverse*
Benedetto colui che viene nel nome tuo, Signor. \rep{2}
\endverse
\endsong

% Fonte: LibrettoAccordi.pdf (Canti Coro), H9. Cori alterni resi come ripetizioni.
% Nell'elenco: Santo, Ordinario.
\beginsong{Osanna nelle altezze}[by={RnS}]
\beginchorus
Santo! Santo! \rep{2}
\endchorus
\beginverse*
I cieli e la terra sono pieni della tua gloria.
Osanna! Osanna! Osanna! Osanna nelle altezze! \rep{2}
\endverse
\beginverse*
Benedetto colui che viene nel nome del Signore.
Osanna! Osanna! Osanna! Osanna nelle altezze! \rep{2}
\endverse
\endsong

% Fonte: LibrettoAccordi.pdf (Canti Coro), H3. Testo confermato dall'utente.
% Nell'elenco: Santo, Ordinario. CD: Verbum Panis (2000).
\beginsong{Sanctus}[by={M. Balduzzi (CD Verbum Panis)}]
\beginverse*
Sanctus, sanctus, sanctus, \rep{2}
Dominus Deus Sabaoth, \rep{2}
Pleni sunt caeli et terra gloria tua.
Hosanna in excelsis. \rep{2}
\endverse
\beginverse*
Benedictus qui venit in nomine Domini, \rep{2}
Hosanna in excelsis. \rep{2}
\endverse
\endsong

% Fonte: LibrettoAccordi.pdf (Canti Coro), H8.
% Nell'elenco: Santo, Ordinario. CD: Giorno senza tramonto.
\beginsong{Santo\soloindice{ (Balduzzi, Giorno senza tramonto)}}[by={M. Balduzzi (CD Giorno senza tramonto)}]
\beginchorus
Osanna! Osanna!
Osanna! Osanna!
Osanna!
\endchorus
\beginverse*
Santo è il Signore, Dio dell'universo.
Santo è il Signore, nostro Dio.
Cieli e terra cantano in coro la gloria.
Cieli e terra cantano in coro la gloria sua.
\endverse
\beginverse*
Benedetto è lui che viene nel nome,
benedetto è lui che viene.
Benedetto è lui che viene nel nome,
benedetto è lui che viene nel nome del Signor.
\endverse
\endsong

% Fonte: LibrettoAccordi.pdf (Canti Coro), H6.
% Nell'elenco: Santo, Ordinario. CD: Come fuoco vivo.
\beginsong{Santo\soloindice{ (Come fuoco vivo)}}[by={Gen Rosso, Gen Verde (CD Come fuoco vivo)}]
\beginverse*
Santo! Santo! Santo è il Signore Dio dell'universo.
Santo! Santo! I cieli e la terra sono pieni della tua gloria.
\endverse
\beginchorus
Osanna nell'alto dei cieli! Osanna nell'alto dei cieli!
\endchorus
\beginverse*
Santo! Santo! Santo è il Signore Dio dell'universo.
Santo! Santo! I cieli e la terra sono pieni della tua gloria.
Benedetto colui che viene nel nome del Signore!
\endverse
\beginverse*
Santo! Santo! Santo!
\endverse
\endsong

% Fonte: LibrettoAccordi.pdf (Canti Coro), H1.
% Nell'elenco: Santo, Ordinario.
\beginsong{Santo\soloindice{ (F. Baggio)}}[by={F. Baggio}]
\beginverse*
Santo, santo, santo il Signore,
Dio dell'universo.
I cieli e la terra sono pieni
della tua gloria.
\endverse
\beginchorus
Osanna nell'alto dei cieli!
Osanna nell'alto dei cieli!
\endchorus
\beginverse*
Benedetto colui che viene
nel nome del Signore.
Osanna! Osanna nell'alto dei cieli! \rep{2}
Osanna! Osanna! Osanna nell'alto dei cieli!
\endverse
\endsong

% Fonte: LibrettoAccordi.pdf (Canti Coro), H11.
% Nell'elenco: righe 61 (Ordinario) e 62 (Settimana Santa), stesso brano secondo l'utente. CD: Il mistero pasquale (2011).
\beginsong{Santo\soloindice{ (Gen Verde, Il mistero pasquale)}}[by={Gen Verde (CD Il mistero pasquale)}]
\beginverse*
Santo, Santo, Santo il Signore Dio dell'universo.
I cieli e la terra sono pieni della tua gloria.
\endverse
\beginchorus
Osanna, osanna, osanna nell'alto dei cieli. \rep{2}
\endchorus
\beginverse*
Benedetto colui che viene nel nome del Signore.
\endverse
\beginverse*
Santo, Santo, Santo.
\endverse
\endsong

% Fonte: LibrettoAccordi.pdf (Canti Coro), H4.
% Nell'elenco: Santo, Ordinario.
\beginsong{Santo il Signore}[by={D. Ricci}]
\beginverse*
Santo il Signore Dio dell'universo,
i cieli e la terra son pieni della tua gloria.
Osanna nell'alto dei cieli e benedetto colui che viene
nel nome del Signore. Osanna nell'alto dei cieli.
\endverse
\beginverse*
È santo, santo, santo il Signore della vita.
È santo, santo, santo il Signore della storia,
nel cielo e sulla terra risplende la sua gloria.
Osanna, osanna, osanna nel cielo sconfinato.
Osanna, osanna, osanna nel canto del creato.
E sia benedetto colui che viene in nome di Dio. \rep{3}
\endverse
\endsong

% Fonte: LibrettoAccordi.pdf (Canti Coro), H2.
% Nell'elenco: Santo, Ordinario.
\beginsong{Santo, santo, santo}[by={D. Ricci}]
\beginverse*
Santo! Santo! Santo
è il Signore Dio dell'universo!
Santo! Santo! Santo!
I cieli e la terra sono pieni della tua gloria!
Osanna! Osanna! Osanna nell'alto dei cieli!
Benedetto colui che viene nel nome del Signore.
Osanna! Osanna! Osanna nell'alto dei cieli! \rep{3}
\endverse
\endsong

% Fonte: LibrettoAccordi.pdf (Canti Coro), H7.
% Nell'elenco: Santo, Ordinario. CD: Xaipe.
\beginsong{Santo (Xaipe)}[by={M. Balduzzi (CD Xaipe)}]
\beginchorus
Osanna! Osanna in excelsis! Osanna!
\endchorus
\beginverse*
Sanctus Dominus Deus Sabaoth! \rep{2}
Pleni, pleni sunt coeli et terra \rep{2}
gloria tua.
\endverse
\beginverse*
Benedictus qui venit \rep{2}
in nomine Domini.
\endverse
\endsong

\end{songs}

% ===================== ALLELUIA E LODE =====================
% !TEX root = ../libretto.tex
% Sezione: Alleluia e Lode (acclamazioni al Vangelo, anche quaresimali)
\momento{Alleluia e Lode}
\begin{songs}{titoli}

% Fonte: LibrettoAccordi.pdf (Canti Coro), E5.
% Nell'elenco: Alleluia, Matrimonio.
\beginsong{Alleluia, beati coloro}[by={F. Baggio, F. Buttazzo}]
\beginchorus
Alleluja! Alleluja! Alleluja! Alleluja! Alleluja! Alleluja! \rep{2}
\endchorus
\beginverse
Beati coloro che sanno ascoltare quello che dice lo Sposo alla Chiesa.
Cristo Signore rivela il suo amore per la sua sposa, la Chiesa.
\endverse
\endsong

% Fonte: LibrettoAccordi.pdf (Canti Coro), E4.
% Nell'elenco: Alleluia, Ordinario.
\beginsong{Alleluia, canto per Cristo}[by={E. Costa}]
\beginchorus
Alleluja! Alleluja! Alleluja! Alleluja! Alleluja! \rep{2}
\endchorus
\beginverse
Canto per Cristo che mi libererà, quando verrà nella gloria,
quando la vita con lui rinascerà. Alleluja! Alleluja!
\endverse
\beginverse
Canto per Cristo: in Lui rifiorirà ogni speranza perduta,
ogni creatura con lui risorgerà. Alleluja! Alleluja!
\endverse
\beginverse
Canto per Cristo: un giorno tornerà! Festa per tutti gli amici,
festa di un mondo che più non morirà. Alleluja! Alleluja!
\endverse
\endsong

% Fonte: LibrettoAccordi.pdf (Canti Coro), E2.
% Nell'elenco: Alleluia, Pentecoste; l'utente lo vuole tra gli Alleluia.
\beginsong{Alleluia (Come fuoco vivo)}[by={Gen Rosso, Gen Verde (CD Come fuoco vivo)}]
\beginverse*
Alleluja! Alleluja! Alleluja!
Alleluja! Alleluja! Alleluja! Alleluja!
Ah, alleluja! Ah, alleluja, alleluja!
Alleluja! Alleluja! Alleluja!
Alleluja! Alleluja! Alleluja! Alleluja! Alleluja!
\endverse
\endsong

% Fonte: LibrettoAccordi.pdf (Canti Coro), E12.
% Nell'elenco: Alleluia, Ordinario.
\beginsong{Alleluia\soloindice{ (Gen Verde)}}[by={Gen Verde}]
\beginverse*
Alleluja! Alleluja! Alleluja! Alleluja! \rep{2}
Alleluja! Alleluja! Alleluja! Ah!
Alleluja! Alleluja! Alleluja! Alleluja! Alleluja! Alleluja!
Alleluja! Alleluja! Alleluja! Alleluja! \rep{2}
\endverse
\endsong

% Fonte: LibrettoAccordi.pdf (Canti Coro), E11. Omesse le indicazioni dei due gruppi.
% Nell'elenco: Alleluia, Ordinario.
\beginsong{Alleluia, lode cosmica}[by={S. Puri}]
\beginchorus
Alleluja! Alleluja! Alleluja! Alleluja!
Alleluja! Alleluja! Alleluja! Alleluja! \rep{2}
\endchorus
\beginverse
Lodino il Signor i cieli, lodino il Signor i mari,
gli angeli, i cieli dei cieli: il suo nome è grande e sublime.
Lode, lode al suo nome.
Sole, luna e stelle ardenti, neve, pioggia, nebbia e fuoco
lodino il suo nome in eterno!
Sia lode al Signor! Sia lode al Signor!
Sia lode al Signore!
\endverse
\beginverse
Lodino il Signor le terre, lodino il Signor i monti,
il vento della tempesta che obbedisce alla sua voce.
Giudici, sovrani tutti, giovani, fanciulle, vecchi
lodino il suo nome in eterno!
Sia lode al Signor! Sia lode al Signor!
Sia lode al Signore!
\endverse
\endsong

% Fonte: LibrettoAccordi.pdf (Canti Coro), E3.
% Nell'elenco: Alleluia, Ordinario.
\beginsong{Alleluia\soloindice{ (OMI)}}[by={OMI}]
\beginverse*
Alleluja! Alleluja!
Alleluja! Alleluja! Alleluja! \rep{2}
\endverse
\endsong

% Fonte: LibrettoAccordi.pdf (Canti Coro), E6.
% Nell'elenco: Alleluia, Ordinario.
\beginsong{Alleluia, ta parole}[by={F. Baggio}]
\beginchorus
Alle-alleluja! Alleluja. Alle-alleluja! Alleluja. \rep{2}
\endchorus
\beginverse
È la tua parola, la verità, che dona luce a questa umanità:
rivela all'uomo il cuore di Dio e porta al mondo unità.
\endverse
\endsong

% Fonte: LibrettoAccordi.pdf (Canti Coro), E9.
% Nell'elenco: Alleluia, Ordinario.
\beginsong{Alleluia\soloindice{ (Verbum Panis)}}[by={M. Balduzzi (CD Verbum Panis)}]
\beginchorus
Alleluja! Alleluja! Alleluja! Alleluja! Alleluja! \rep{2}
\endchorus
\beginverse*
Alleluja! Alleluja! \rep{4}
\endverse
\endsong

% Fonte: LibrettoAccordi.pdf (Canti Coro), E7.
% Nell'elenco: Alleluia (acclamazione), Quaresima. Spostato in "Alleluia e Lode" su indicazione dell'utente.
\beginsong{Lode a te, o Cristo\soloindice{ (F. Buttazzo)}}[by={F. Buttazzo}]
\beginchorus
Lode a Te, o Cristo, Verbo di salvezza!
Lode a Te, o Cristo, Parola di Dio per noi! \rep{2}
\endchorus
\beginverse
Come la luce illumina ogni cosa,
la tua Parola rivela a noi il tuo volto.
Come la pioggia scende nella terra,
la tua Parola entra nel nostro cuore.
\endverse
\beginverse
Seme da te piantato in terra buona,
la tua Parola in noi porterà i suoi frutti.
Pane che tu ci doni in ogni giorno,
la tua Parola è cibo per noi di vita.
\endverse
\endsong

% Fonte: LibrettoAccordi.pdf (Canti Coro), E13.
% Nell'elenco: Alleluia (acclamazione), Quaresima. Spostato in "Alleluia e Lode" su indicazione dell'utente.
\beginsong{Lode a te, o Cristo\soloindice{ (Gen Verde)}}[by={Gen Verde}]
\beginchorus
Lode a Te, o Cristo,
Re di eterna gloria. \rep{2}
\endchorus
\beginverse
Signore Tu sei veramente
il Salvatore del mondo.
Dammi dell'acqua viva
perché non abbia più sete.
\endverse
\beginverse
Chi beve di quest'acqua
avrà di nuovo sete.
Ma chi beve dell'acqua che io gli darò
non avrà mai più sete.
\endverse
\endsong

% Fonte: LibrettoAccordi.pdf (Canti Coro), E10.
% Nell'elenco: Alleluia (acclamazione), Quaresima. Spostato in "Alleluia e Lode" su indicazione dell'utente.
\beginsong{Lode e gloria}[by={P. Spoladore}]
\beginchorus
Lode e gloria a Te, o Signore! \rep{2}
\endchorus
\beginverse
Date lode al Signore o figli di Dio,
benedite il suo nome su tutta la terra.
Il suo amore per noi è fedele per sempre,
il suo amore per noi è per sempre.
\endverse
\beginverse
Acclamate al Signore da tutta la terra,
ci ha creati e voluti, a lui apparteniamo.
Il Signore ci guida e ci tiene per mano,
senza fine è il suo amore per noi.
\endverse
\endsong

\end{songs}
 in libretto.tex (o ricopiare le due parti in canti/).

% ===================== SANTO =====================
% !TEX root = ../libretto.tex
% Sezione: Santo
% Nei brani a due cori (Uomini/Donne) le parti sono state unite in un solo testo.
\momento{Santo}
\begin{songs}{titoli}

% Fonte: LibrettoAccordi.pdf (Canti Coro), H10.
% Nell'elenco: Santo, Ordinario.
\beginsong{Osanna eh}[by={RnS}]
\beginchorus
Santo! Santo! Osanna! \rep{2}
Osanna eh! Osanna eh! Osanna a Cristo Signor! \rep{2}
\endchorus
\beginverse*
I cieli e la terra, o Signore, sono pieni di te. \rep{2}
\endverse
\beginverse*
Benedetto colui che viene nel nome tuo, Signor. \rep{2}
\endverse
\endsong

% Fonte: LibrettoAccordi.pdf (Canti Coro), H9. Cori alterni resi come ripetizioni.
% Nell'elenco: Santo, Ordinario.
\beginsong{Osanna nelle altezze}[by={RnS}]
\beginchorus
Santo! Santo! \rep{2}
\endchorus
\beginverse*
I cieli e la terra sono pieni della tua gloria.
Osanna! Osanna! Osanna! Osanna nelle altezze! \rep{2}
\endverse
\beginverse*
Benedetto colui che viene nel nome del Signore.
Osanna! Osanna! Osanna! Osanna nelle altezze! \rep{2}
\endverse
\endsong

% Fonte: LibrettoAccordi.pdf (Canti Coro), H3. Testo confermato dall'utente.
% Nell'elenco: Santo, Ordinario. CD: Verbum Panis (2000).
\beginsong{Sanctus}[by={M. Balduzzi (CD Verbum Panis)}]
\beginverse*
Sanctus, sanctus, sanctus, \rep{2}
Dominus Deus Sabaoth, \rep{2}
Pleni sunt caeli et terra gloria tua.
Hosanna in excelsis. \rep{2}
\endverse
\beginverse*
Benedictus qui venit in nomine Domini, \rep{2}
Hosanna in excelsis. \rep{2}
\endverse
\endsong

% Fonte: LibrettoAccordi.pdf (Canti Coro), H8.
% Nell'elenco: Santo, Ordinario. CD: Giorno senza tramonto.
\beginsong{Santo\soloindice{ (Balduzzi, Giorno senza tramonto)}}[by={M. Balduzzi (CD Giorno senza tramonto)}]
\beginchorus
Osanna! Osanna!
Osanna! Osanna!
Osanna!
\endchorus
\beginverse*
Santo è il Signore, Dio dell'universo.
Santo è il Signore, nostro Dio.
Cieli e terra cantano in coro la gloria.
Cieli e terra cantano in coro la gloria sua.
\endverse
\beginverse*
Benedetto è lui che viene nel nome,
benedetto è lui che viene.
Benedetto è lui che viene nel nome,
benedetto è lui che viene nel nome del Signor.
\endverse
\endsong

% Fonte: LibrettoAccordi.pdf (Canti Coro), H6.
% Nell'elenco: Santo, Ordinario. CD: Come fuoco vivo.
\beginsong{Santo\soloindice{ (Come fuoco vivo)}}[by={Gen Rosso, Gen Verde (CD Come fuoco vivo)}]
\beginverse*
Santo! Santo! Santo è il Signore Dio dell'universo.
Santo! Santo! I cieli e la terra sono pieni della tua gloria.
\endverse
\beginchorus
Osanna nell'alto dei cieli! Osanna nell'alto dei cieli!
\endchorus
\beginverse*
Santo! Santo! Santo è il Signore Dio dell'universo.
Santo! Santo! I cieli e la terra sono pieni della tua gloria.
Benedetto colui che viene nel nome del Signore!
\endverse
\beginverse*
Santo! Santo! Santo!
\endverse
\endsong

% Fonte: LibrettoAccordi.pdf (Canti Coro), H1.
% Nell'elenco: Santo, Ordinario.
\beginsong{Santo\soloindice{ (F. Baggio)}}[by={F. Baggio}]
\beginverse*
Santo, santo, santo il Signore,
Dio dell'universo.
I cieli e la terra sono pieni
della tua gloria.
\endverse
\beginchorus
Osanna nell'alto dei cieli!
Osanna nell'alto dei cieli!
\endchorus
\beginverse*
Benedetto colui che viene
nel nome del Signore.
Osanna! Osanna nell'alto dei cieli! \rep{2}
Osanna! Osanna! Osanna nell'alto dei cieli!
\endverse
\endsong

% Fonte: LibrettoAccordi.pdf (Canti Coro), H11.
% Nell'elenco: righe 61 (Ordinario) e 62 (Settimana Santa), stesso brano secondo l'utente. CD: Il mistero pasquale (2011).
\beginsong{Santo\soloindice{ (Gen Verde, Il mistero pasquale)}}[by={Gen Verde (CD Il mistero pasquale)}]
\beginverse*
Santo, Santo, Santo il Signore Dio dell'universo.
I cieli e la terra sono pieni della tua gloria.
\endverse
\beginchorus
Osanna, osanna, osanna nell'alto dei cieli. \rep{2}
\endchorus
\beginverse*
Benedetto colui che viene nel nome del Signore.
\endverse
\beginverse*
Santo, Santo, Santo.
\endverse
\endsong

% Fonte: LibrettoAccordi.pdf (Canti Coro), H4.
% Nell'elenco: Santo, Ordinario.
\beginsong{Santo il Signore}[by={D. Ricci}]
\beginverse*
Santo il Signore Dio dell'universo,
i cieli e la terra son pieni della tua gloria.
Osanna nell'alto dei cieli e benedetto colui che viene
nel nome del Signore. Osanna nell'alto dei cieli.
\endverse
\beginverse*
È santo, santo, santo il Signore della vita.
È santo, santo, santo il Signore della storia,
nel cielo e sulla terra risplende la sua gloria.
Osanna, osanna, osanna nel cielo sconfinato.
Osanna, osanna, osanna nel canto del creato.
E sia benedetto colui che viene in nome di Dio. \rep{3}
\endverse
\endsong

% Fonte: LibrettoAccordi.pdf (Canti Coro), H2.
% Nell'elenco: Santo, Ordinario.
\beginsong{Santo, santo, santo}[by={D. Ricci}]
\beginverse*
Santo! Santo! Santo
è il Signore Dio dell'universo!
Santo! Santo! Santo!
I cieli e la terra sono pieni della tua gloria!
Osanna! Osanna! Osanna nell'alto dei cieli!
Benedetto colui che viene nel nome del Signore.
Osanna! Osanna! Osanna nell'alto dei cieli! \rep{3}
\endverse
\endsong

% Fonte: LibrettoAccordi.pdf (Canti Coro), H7.
% Nell'elenco: Santo, Ordinario. CD: Xaipe.
\beginsong{Santo (Xaipe)}[by={M. Balduzzi (CD Xaipe)}]
\beginchorus
Osanna! Osanna in excelsis! Osanna!
\endchorus
\beginverse*
Sanctus Dominus Deus Sabaoth! \rep{2}
Pleni, pleni sunt coeli et terra \rep{2}
gloria tua.
\endverse
\beginverse*
Benedictus qui venit \rep{2}
in nomine Domini.
\endverse
\endsong

\end{songs}

% ===================== ALLELUIA E LODE =====================
% !TEX root = ../libretto.tex
% Sezione: Alleluia e Lode (acclamazioni al Vangelo, anche quaresimali)
\momento{Alleluia e Lode}
\begin{songs}{titoli}

% Fonte: LibrettoAccordi.pdf (Canti Coro), E5.
% Nell'elenco: Alleluia, Matrimonio.
\beginsong{Alleluia, beati coloro}[by={F. Baggio, F. Buttazzo}]
\beginchorus
Alleluja! Alleluja! Alleluja! Alleluja! Alleluja! Alleluja! \rep{2}
\endchorus
\beginverse
Beati coloro che sanno ascoltare quello che dice lo Sposo alla Chiesa.
Cristo Signore rivela il suo amore per la sua sposa, la Chiesa.
\endverse
\endsong

% Fonte: LibrettoAccordi.pdf (Canti Coro), E4.
% Nell'elenco: Alleluia, Ordinario.
\beginsong{Alleluia, canto per Cristo}[by={E. Costa}]
\beginchorus
Alleluja! Alleluja! Alleluja! Alleluja! Alleluja! \rep{2}
\endchorus
\beginverse
Canto per Cristo che mi libererà, quando verrà nella gloria,
quando la vita con lui rinascerà. Alleluja! Alleluja!
\endverse
\beginverse
Canto per Cristo: in Lui rifiorirà ogni speranza perduta,
ogni creatura con lui risorgerà. Alleluja! Alleluja!
\endverse
\beginverse
Canto per Cristo: un giorno tornerà! Festa per tutti gli amici,
festa di un mondo che più non morirà. Alleluja! Alleluja!
\endverse
\endsong

% Fonte: LibrettoAccordi.pdf (Canti Coro), E2.
% Nell'elenco: Alleluia, Pentecoste; l'utente lo vuole tra gli Alleluia.
\beginsong{Alleluia (Come fuoco vivo)}[by={Gen Rosso, Gen Verde (CD Come fuoco vivo)}]
\beginverse*
Alleluja! Alleluja! Alleluja!
Alleluja! Alleluja! Alleluja! Alleluja!
Ah, alleluja! Ah, alleluja, alleluja!
Alleluja! Alleluja! Alleluja!
Alleluja! Alleluja! Alleluja! Alleluja! Alleluja!
\endverse
\endsong

% Fonte: LibrettoAccordi.pdf (Canti Coro), E12.
% Nell'elenco: Alleluia, Ordinario.
\beginsong{Alleluia\soloindice{ (Gen Verde)}}[by={Gen Verde}]
\beginverse*
Alleluja! Alleluja! Alleluja! Alleluja! \rep{2}
Alleluja! Alleluja! Alleluja! Ah!
Alleluja! Alleluja! Alleluja! Alleluja! Alleluja! Alleluja!
Alleluja! Alleluja! Alleluja! Alleluja! \rep{2}
\endverse
\endsong

% Fonte: LibrettoAccordi.pdf (Canti Coro), E11. Omesse le indicazioni dei due gruppi.
% Nell'elenco: Alleluia, Ordinario.
\beginsong{Alleluia, lode cosmica}[by={S. Puri}]
\beginchorus
Alleluja! Alleluja! Alleluja! Alleluja!
Alleluja! Alleluja! Alleluja! Alleluja! \rep{2}
\endchorus
\beginverse
Lodino il Signor i cieli, lodino il Signor i mari,
gli angeli, i cieli dei cieli: il suo nome è grande e sublime.
Lode, lode al suo nome.
Sole, luna e stelle ardenti, neve, pioggia, nebbia e fuoco
lodino il suo nome in eterno!
Sia lode al Signor! Sia lode al Signor!
Sia lode al Signore!
\endverse
\beginverse
Lodino il Signor le terre, lodino il Signor i monti,
il vento della tempesta che obbedisce alla sua voce.
Giudici, sovrani tutti, giovani, fanciulle, vecchi
lodino il suo nome in eterno!
Sia lode al Signor! Sia lode al Signor!
Sia lode al Signore!
\endverse
\endsong

% Fonte: LibrettoAccordi.pdf (Canti Coro), E3.
% Nell'elenco: Alleluia, Ordinario.
\beginsong{Alleluia\soloindice{ (OMI)}}[by={OMI}]
\beginverse*
Alleluja! Alleluja!
Alleluja! Alleluja! Alleluja! \rep{2}
\endverse
\endsong

% Fonte: LibrettoAccordi.pdf (Canti Coro), E6.
% Nell'elenco: Alleluia, Ordinario.
\beginsong{Alleluia, ta parole}[by={F. Baggio}]
\beginchorus
Alle-alleluja! Alleluja. Alle-alleluja! Alleluja. \rep{2}
\endchorus
\beginverse
È la tua parola, la verità, che dona luce a questa umanità:
rivela all'uomo il cuore di Dio e porta al mondo unità.
\endverse
\endsong

% Fonte: LibrettoAccordi.pdf (Canti Coro), E9.
% Nell'elenco: Alleluia, Ordinario.
\beginsong{Alleluia\soloindice{ (Verbum Panis)}}[by={M. Balduzzi (CD Verbum Panis)}]
\beginchorus
Alleluja! Alleluja! Alleluja! Alleluja! Alleluja! \rep{2}
\endchorus
\beginverse*
Alleluja! Alleluja! \rep{4}
\endverse
\endsong

% Fonte: LibrettoAccordi.pdf (Canti Coro), E7.
% Nell'elenco: Alleluia (acclamazione), Quaresima. Spostato in "Alleluia e Lode" su indicazione dell'utente.
\beginsong{Lode a te, o Cristo\soloindice{ (F. Buttazzo)}}[by={F. Buttazzo}]
\beginchorus
Lode a Te, o Cristo, Verbo di salvezza!
Lode a Te, o Cristo, Parola di Dio per noi! \rep{2}
\endchorus
\beginverse
Come la luce illumina ogni cosa,
la tua Parola rivela a noi il tuo volto.
Come la pioggia scende nella terra,
la tua Parola entra nel nostro cuore.
\endverse
\beginverse
Seme da te piantato in terra buona,
la tua Parola in noi porterà i suoi frutti.
Pane che tu ci doni in ogni giorno,
la tua Parola è cibo per noi di vita.
\endverse
\endsong

% Fonte: LibrettoAccordi.pdf (Canti Coro), E13.
% Nell'elenco: Alleluia (acclamazione), Quaresima. Spostato in "Alleluia e Lode" su indicazione dell'utente.
\beginsong{Lode a te, o Cristo\soloindice{ (Gen Verde)}}[by={Gen Verde}]
\beginchorus
Lode a Te, o Cristo,
Re di eterna gloria. \rep{2}
\endchorus
\beginverse
Signore Tu sei veramente
il Salvatore del mondo.
Dammi dell'acqua viva
perché non abbia più sete.
\endverse
\beginverse
Chi beve di quest'acqua
avrà di nuovo sete.
Ma chi beve dell'acqua che io gli darò
non avrà mai più sete.
\endverse
\endsong

% Fonte: LibrettoAccordi.pdf (Canti Coro), E10.
% Nell'elenco: Alleluia (acclamazione), Quaresima. Spostato in "Alleluia e Lode" su indicazione dell'utente.
\beginsong{Lode e gloria}[by={P. Spoladore}]
\beginchorus
Lode e gloria a Te, o Signore! \rep{2}
\endchorus
\beginverse
Date lode al Signore o figli di Dio,
benedite il suo nome su tutta la terra.
Il suo amore per noi è fedele per sempre,
il suo amore per noi è per sempre.
\endverse
\beginverse
Acclamate al Signore da tutta la terra,
ci ha creati e voluti, a lui apparteniamo.
Il Signore ci guida e ci tiene per mano,
senza fine è il suo amore per noi.
\endverse
\endsong

\end{songs}
 in libretto.tex (o ricopiare le due parti in canti/).

% ===================== SANTO =====================
% !TEX root = ../libretto.tex
% Sezione: Santo
% Nei brani a due cori (Uomini/Donne) le parti sono state unite in un solo testo.
\momento{Santo}
\begin{songs}{titoli}

% Fonte: LibrettoAccordi.pdf (Canti Coro), H10.
% Nell'elenco: Santo, Ordinario.
\beginsong{Osanna eh}[by={RnS}]
\beginchorus
Santo! Santo! Osanna! \rep{2}
Osanna eh! Osanna eh! Osanna a Cristo Signor! \rep{2}
\endchorus
\beginverse*
I cieli e la terra, o Signore, sono pieni di te. \rep{2}
\endverse
\beginverse*
Benedetto colui che viene nel nome tuo, Signor. \rep{2}
\endverse
\endsong

% Fonte: LibrettoAccordi.pdf (Canti Coro), H9. Cori alterni resi come ripetizioni.
% Nell'elenco: Santo, Ordinario.
\beginsong{Osanna nelle altezze}[by={RnS}]
\beginchorus
Santo! Santo! \rep{2}
\endchorus
\beginverse*
I cieli e la terra sono pieni della tua gloria.
Osanna! Osanna! Osanna! Osanna nelle altezze! \rep{2}
\endverse
\beginverse*
Benedetto colui che viene nel nome del Signore.
Osanna! Osanna! Osanna! Osanna nelle altezze! \rep{2}
\endverse
\endsong

% Fonte: LibrettoAccordi.pdf (Canti Coro), H3. Testo confermato dall'utente.
% Nell'elenco: Santo, Ordinario. CD: Verbum Panis (2000).
\beginsong{Sanctus}[by={M. Balduzzi (CD Verbum Panis)}]
\beginverse*
Sanctus, sanctus, sanctus, \rep{2}
Dominus Deus Sabaoth, \rep{2}
Pleni sunt caeli et terra gloria tua.
Hosanna in excelsis. \rep{2}
\endverse
\beginverse*
Benedictus qui venit in nomine Domini, \rep{2}
Hosanna in excelsis. \rep{2}
\endverse
\endsong

% Fonte: LibrettoAccordi.pdf (Canti Coro), H8.
% Nell'elenco: Santo, Ordinario. CD: Giorno senza tramonto.
\beginsong{Santo\soloindice{ (Balduzzi, Giorno senza tramonto)}}[by={M. Balduzzi (CD Giorno senza tramonto)}]
\beginchorus
Osanna! Osanna!
Osanna! Osanna!
Osanna!
\endchorus
\beginverse*
Santo è il Signore, Dio dell'universo.
Santo è il Signore, nostro Dio.
Cieli e terra cantano in coro la gloria.
Cieli e terra cantano in coro la gloria sua.
\endverse
\beginverse*
Benedetto è lui che viene nel nome,
benedetto è lui che viene.
Benedetto è lui che viene nel nome,
benedetto è lui che viene nel nome del Signor.
\endverse
\endsong

% Fonte: LibrettoAccordi.pdf (Canti Coro), H6.
% Nell'elenco: Santo, Ordinario. CD: Come fuoco vivo.
\beginsong{Santo\soloindice{ (Come fuoco vivo)}}[by={Gen Rosso, Gen Verde (CD Come fuoco vivo)}]
\beginverse*
Santo! Santo! Santo è il Signore Dio dell'universo.
Santo! Santo! I cieli e la terra sono pieni della tua gloria.
\endverse
\beginchorus
Osanna nell'alto dei cieli! Osanna nell'alto dei cieli!
\endchorus
\beginverse*
Santo! Santo! Santo è il Signore Dio dell'universo.
Santo! Santo! I cieli e la terra sono pieni della tua gloria.
Benedetto colui che viene nel nome del Signore!
\endverse
\beginverse*
Santo! Santo! Santo!
\endverse
\endsong

% Fonte: LibrettoAccordi.pdf (Canti Coro), H1.
% Nell'elenco: Santo, Ordinario.
\beginsong{Santo\soloindice{ (F. Baggio)}}[by={F. Baggio}]
\beginverse*
Santo, santo, santo il Signore,
Dio dell'universo.
I cieli e la terra sono pieni
della tua gloria.
\endverse
\beginchorus
Osanna nell'alto dei cieli!
Osanna nell'alto dei cieli!
\endchorus
\beginverse*
Benedetto colui che viene
nel nome del Signore.
Osanna! Osanna nell'alto dei cieli! \rep{2}
Osanna! Osanna! Osanna nell'alto dei cieli!
\endverse
\endsong

% Fonte: LibrettoAccordi.pdf (Canti Coro), H11.
% Nell'elenco: righe 61 (Ordinario) e 62 (Settimana Santa), stesso brano secondo l'utente. CD: Il mistero pasquale (2011).
\beginsong{Santo\soloindice{ (Gen Verde, Il mistero pasquale)}}[by={Gen Verde (CD Il mistero pasquale)}]
\beginverse*
Santo, Santo, Santo il Signore Dio dell'universo.
I cieli e la terra sono pieni della tua gloria.
\endverse
\beginchorus
Osanna, osanna, osanna nell'alto dei cieli. \rep{2}
\endchorus
\beginverse*
Benedetto colui che viene nel nome del Signore.
\endverse
\beginverse*
Santo, Santo, Santo.
\endverse
\endsong

% Fonte: LibrettoAccordi.pdf (Canti Coro), H4.
% Nell'elenco: Santo, Ordinario.
\beginsong{Santo il Signore}[by={D. Ricci}]
\beginverse*
Santo il Signore Dio dell'universo,
i cieli e la terra son pieni della tua gloria.
Osanna nell'alto dei cieli e benedetto colui che viene
nel nome del Signore. Osanna nell'alto dei cieli.
\endverse
\beginverse*
È santo, santo, santo il Signore della vita.
È santo, santo, santo il Signore della storia,
nel cielo e sulla terra risplende la sua gloria.
Osanna, osanna, osanna nel cielo sconfinato.
Osanna, osanna, osanna nel canto del creato.
E sia benedetto colui che viene in nome di Dio. \rep{3}
\endverse
\endsong

% Fonte: LibrettoAccordi.pdf (Canti Coro), H2.
% Nell'elenco: Santo, Ordinario.
\beginsong{Santo, santo, santo}[by={D. Ricci}]
\beginverse*
Santo! Santo! Santo
è il Signore Dio dell'universo!
Santo! Santo! Santo!
I cieli e la terra sono pieni della tua gloria!
Osanna! Osanna! Osanna nell'alto dei cieli!
Benedetto colui che viene nel nome del Signore.
Osanna! Osanna! Osanna nell'alto dei cieli! \rep{3}
\endverse
\endsong

% Fonte: LibrettoAccordi.pdf (Canti Coro), H7.
% Nell'elenco: Santo, Ordinario. CD: Xaipe.
\beginsong{Santo (Xaipe)}[by={M. Balduzzi (CD Xaipe)}]
\beginchorus
Osanna! Osanna in excelsis! Osanna!
\endchorus
\beginverse*
Sanctus Dominus Deus Sabaoth! \rep{2}
Pleni, pleni sunt coeli et terra \rep{2}
gloria tua.
\endverse
\beginverse*
Benedictus qui venit \rep{2}
in nomine Domini.
\endverse
\endsong

\end{songs}

% ===================== ALLELUIA E LODE =====================
% !TEX root = ../libretto.tex
% Sezione: Alleluia e Lode (acclamazioni al Vangelo, anche quaresimali)
\momento{Alleluia e Lode}
\begin{songs}{titoli}

% Fonte: LibrettoAccordi.pdf (Canti Coro), E5.
% Nell'elenco: Alleluia, Matrimonio.
\beginsong{Alleluia, beati coloro}[by={F. Baggio, F. Buttazzo}]
\beginchorus
Alleluja! Alleluja! Alleluja! Alleluja! Alleluja! Alleluja! \rep{2}
\endchorus
\beginverse
Beati coloro che sanno ascoltare quello che dice lo Sposo alla Chiesa.
Cristo Signore rivela il suo amore per la sua sposa, la Chiesa.
\endverse
\endsong

% Fonte: LibrettoAccordi.pdf (Canti Coro), E4.
% Nell'elenco: Alleluia, Ordinario.
\beginsong{Alleluia, canto per Cristo}[by={E. Costa}]
\beginchorus
Alleluja! Alleluja! Alleluja! Alleluja! Alleluja! \rep{2}
\endchorus
\beginverse
Canto per Cristo che mi libererà, quando verrà nella gloria,
quando la vita con lui rinascerà. Alleluja! Alleluja!
\endverse
\beginverse
Canto per Cristo: in Lui rifiorirà ogni speranza perduta,
ogni creatura con lui risorgerà. Alleluja! Alleluja!
\endverse
\beginverse
Canto per Cristo: un giorno tornerà! Festa per tutti gli amici,
festa di un mondo che più non morirà. Alleluja! Alleluja!
\endverse
\endsong

% Fonte: LibrettoAccordi.pdf (Canti Coro), E2.
% Nell'elenco: Alleluia, Pentecoste; l'utente lo vuole tra gli Alleluia.
\beginsong{Alleluia (Come fuoco vivo)}[by={Gen Rosso, Gen Verde (CD Come fuoco vivo)}]
\beginverse*
Alleluja! Alleluja! Alleluja!
Alleluja! Alleluja! Alleluja! Alleluja!
Ah, alleluja! Ah, alleluja, alleluja!
Alleluja! Alleluja! Alleluja!
Alleluja! Alleluja! Alleluja! Alleluja! Alleluja!
\endverse
\endsong

% Fonte: LibrettoAccordi.pdf (Canti Coro), E12.
% Nell'elenco: Alleluia, Ordinario.
\beginsong{Alleluia\soloindice{ (Gen Verde)}}[by={Gen Verde}]
\beginverse*
Alleluja! Alleluja! Alleluja! Alleluja! \rep{2}
Alleluja! Alleluja! Alleluja! Ah!
Alleluja! Alleluja! Alleluja! Alleluja! Alleluja! Alleluja!
Alleluja! Alleluja! Alleluja! Alleluja! \rep{2}
\endverse
\endsong

% Fonte: LibrettoAccordi.pdf (Canti Coro), E11. Omesse le indicazioni dei due gruppi.
% Nell'elenco: Alleluia, Ordinario.
\beginsong{Alleluia, lode cosmica}[by={S. Puri}]
\beginchorus
Alleluja! Alleluja! Alleluja! Alleluja!
Alleluja! Alleluja! Alleluja! Alleluja! \rep{2}
\endchorus
\beginverse
Lodino il Signor i cieli, lodino il Signor i mari,
gli angeli, i cieli dei cieli: il suo nome è grande e sublime.
Lode, lode al suo nome.
Sole, luna e stelle ardenti, neve, pioggia, nebbia e fuoco
lodino il suo nome in eterno!
Sia lode al Signor! Sia lode al Signor!
Sia lode al Signore!
\endverse
\beginverse
Lodino il Signor le terre, lodino il Signor i monti,
il vento della tempesta che obbedisce alla sua voce.
Giudici, sovrani tutti, giovani, fanciulle, vecchi
lodino il suo nome in eterno!
Sia lode al Signor! Sia lode al Signor!
Sia lode al Signore!
\endverse
\endsong

% Fonte: LibrettoAccordi.pdf (Canti Coro), E3.
% Nell'elenco: Alleluia, Ordinario.
\beginsong{Alleluia\soloindice{ (OMI)}}[by={OMI}]
\beginverse*
Alleluja! Alleluja!
Alleluja! Alleluja! Alleluja! \rep{2}
\endverse
\endsong

% Fonte: LibrettoAccordi.pdf (Canti Coro), E6.
% Nell'elenco: Alleluia, Ordinario.
\beginsong{Alleluia, ta parole}[by={F. Baggio}]
\beginchorus
Alle-alleluja! Alleluja. Alle-alleluja! Alleluja. \rep{2}
\endchorus
\beginverse
È la tua parola, la verità, che dona luce a questa umanità:
rivela all'uomo il cuore di Dio e porta al mondo unità.
\endverse
\endsong

% Fonte: LibrettoAccordi.pdf (Canti Coro), E9.
% Nell'elenco: Alleluia, Ordinario.
\beginsong{Alleluia\soloindice{ (Verbum Panis)}}[by={M. Balduzzi (CD Verbum Panis)}]
\beginchorus
Alleluja! Alleluja! Alleluja! Alleluja! Alleluja! \rep{2}
\endchorus
\beginverse*
Alleluja! Alleluja! \rep{4}
\endverse
\endsong

% Fonte: LibrettoAccordi.pdf (Canti Coro), E7.
% Nell'elenco: Alleluia (acclamazione), Quaresima. Spostato in "Alleluia e Lode" su indicazione dell'utente.
\beginsong{Lode a te, o Cristo\soloindice{ (F. Buttazzo)}}[by={F. Buttazzo}]
\beginchorus
Lode a Te, o Cristo, Verbo di salvezza!
Lode a Te, o Cristo, Parola di Dio per noi! \rep{2}
\endchorus
\beginverse
Come la luce illumina ogni cosa,
la tua Parola rivela a noi il tuo volto.
Come la pioggia scende nella terra,
la tua Parola entra nel nostro cuore.
\endverse
\beginverse
Seme da te piantato in terra buona,
la tua Parola in noi porterà i suoi frutti.
Pane che tu ci doni in ogni giorno,
la tua Parola è cibo per noi di vita.
\endverse
\endsong

% Fonte: LibrettoAccordi.pdf (Canti Coro), E13.
% Nell'elenco: Alleluia (acclamazione), Quaresima. Spostato in "Alleluia e Lode" su indicazione dell'utente.
\beginsong{Lode a te, o Cristo\soloindice{ (Gen Verde)}}[by={Gen Verde}]
\beginchorus
Lode a Te, o Cristo,
Re di eterna gloria. \rep{2}
\endchorus
\beginverse
Signore Tu sei veramente
il Salvatore del mondo.
Dammi dell'acqua viva
perché non abbia più sete.
\endverse
\beginverse
Chi beve di quest'acqua
avrà di nuovo sete.
Ma chi beve dell'acqua che io gli darò
non avrà mai più sete.
\endverse
\endsong

% Fonte: LibrettoAccordi.pdf (Canti Coro), E10.
% Nell'elenco: Alleluia (acclamazione), Quaresima. Spostato in "Alleluia e Lode" su indicazione dell'utente.
\beginsong{Lode e gloria}[by={P. Spoladore}]
\beginchorus
Lode e gloria a Te, o Signore! \rep{2}
\endchorus
\beginverse
Date lode al Signore o figli di Dio,
benedite il suo nome su tutta la terra.
Il suo amore per noi è fedele per sempre,
il suo amore per noi è per sempre.
\endverse
\beginverse
Acclamate al Signore da tutta la terra,
ci ha creati e voluti, a lui apparteniamo.
Il Signore ci guida e ci tiene per mano,
senza fine è il suo amore per noi.
\endverse
\endsong

\end{songs}
 in libretto.tex (o ricopiare le due parti in canti/).

% ===================== SANTO =====================
% !TEX root = ../libretto.tex
% Sezione: Santo
% Nei brani a due cori (Uomini/Donne) le parti sono state unite in un solo testo.
\momento{Santo}
\begin{songs}{titoli}

% Fonte: LibrettoAccordi.pdf (Canti Coro), H10.
% Nell'elenco: Santo, Ordinario.
\beginsong{Osanna eh}[by={RnS}]
\beginchorus
Santo! Santo! Osanna! \rep{2}
Osanna eh! Osanna eh! Osanna a Cristo Signor! \rep{2}
\endchorus
\beginverse*
I cieli e la terra, o Signore, sono pieni di te. \rep{2}
\endverse
\beginverse*
Benedetto colui che viene nel nome tuo, Signor. \rep{2}
\endverse
\endsong

% Fonte: LibrettoAccordi.pdf (Canti Coro), H9. Cori alterni resi come ripetizioni.
% Nell'elenco: Santo, Ordinario.
\beginsong{Osanna nelle altezze}[by={RnS}]
\beginchorus
Santo! Santo! \rep{2}
\endchorus
\beginverse*
I cieli e la terra sono pieni della tua gloria.
Osanna! Osanna! Osanna! Osanna nelle altezze! \rep{2}
\endverse
\beginverse*
Benedetto colui che viene nel nome del Signore.
Osanna! Osanna! Osanna! Osanna nelle altezze! \rep{2}
\endverse
\endsong

% Fonte: LibrettoAccordi.pdf (Canti Coro), H3. Testo confermato dall'utente.
% Nell'elenco: Santo, Ordinario. CD: Verbum Panis (2000).
\beginsong{Sanctus}[by={M. Balduzzi (CD Verbum Panis)}]
\beginverse*
Sanctus, sanctus, sanctus, \rep{2}
Dominus Deus Sabaoth, \rep{2}
Pleni sunt caeli et terra gloria tua.
Hosanna in excelsis. \rep{2}
\endverse
\beginverse*
Benedictus qui venit in nomine Domini, \rep{2}
Hosanna in excelsis. \rep{2}
\endverse
\endsong

% Fonte: LibrettoAccordi.pdf (Canti Coro), H8.
% Nell'elenco: Santo, Ordinario. CD: Giorno senza tramonto.
\beginsong{Santo\soloindice{ (Balduzzi, Giorno senza tramonto)}}[by={M. Balduzzi (CD Giorno senza tramonto)}]
\beginchorus
Osanna! Osanna!
Osanna! Osanna!
Osanna!
\endchorus
\beginverse*
Santo è il Signore, Dio dell'universo.
Santo è il Signore, nostro Dio.
Cieli e terra cantano in coro la gloria.
Cieli e terra cantano in coro la gloria sua.
\endverse
\beginverse*
Benedetto è lui che viene nel nome,
benedetto è lui che viene.
Benedetto è lui che viene nel nome,
benedetto è lui che viene nel nome del Signor.
\endverse
\endsong

% Fonte: LibrettoAccordi.pdf (Canti Coro), H6.
% Nell'elenco: Santo, Ordinario. CD: Come fuoco vivo.
\beginsong{Santo\soloindice{ (Come fuoco vivo)}}[by={Gen Rosso, Gen Verde (CD Come fuoco vivo)}]
\beginverse*
Santo! Santo! Santo è il Signore Dio dell'universo.
Santo! Santo! I cieli e la terra sono pieni della tua gloria.
\endverse
\beginchorus
Osanna nell'alto dei cieli! Osanna nell'alto dei cieli!
\endchorus
\beginverse*
Santo! Santo! Santo è il Signore Dio dell'universo.
Santo! Santo! I cieli e la terra sono pieni della tua gloria.
Benedetto colui che viene nel nome del Signore!
\endverse
\beginverse*
Santo! Santo! Santo!
\endverse
\endsong

% Fonte: LibrettoAccordi.pdf (Canti Coro), H1.
% Nell'elenco: Santo, Ordinario.
\beginsong{Santo\soloindice{ (F. Baggio)}}[by={F. Baggio}]
\beginverse*
Santo, santo, santo il Signore,
Dio dell'universo.
I cieli e la terra sono pieni
della tua gloria.
\endverse
\beginchorus
Osanna nell'alto dei cieli!
Osanna nell'alto dei cieli!
\endchorus
\beginverse*
Benedetto colui che viene
nel nome del Signore.
Osanna! Osanna nell'alto dei cieli! \rep{2}
Osanna! Osanna! Osanna nell'alto dei cieli!
\endverse
\endsong

% Fonte: LibrettoAccordi.pdf (Canti Coro), H11.
% Nell'elenco: righe 61 (Ordinario) e 62 (Settimana Santa), stesso brano secondo l'utente. CD: Il mistero pasquale (2011).
\beginsong{Santo\soloindice{ (Gen Verde, Il mistero pasquale)}}[by={Gen Verde (CD Il mistero pasquale)}]
\beginverse*
Santo, Santo, Santo il Signore Dio dell'universo.
I cieli e la terra sono pieni della tua gloria.
\endverse
\beginchorus
Osanna, osanna, osanna nell'alto dei cieli. \rep{2}
\endchorus
\beginverse*
Benedetto colui che viene nel nome del Signore.
\endverse
\beginverse*
Santo, Santo, Santo.
\endverse
\endsong

% Fonte: LibrettoAccordi.pdf (Canti Coro), H4.
% Nell'elenco: Santo, Ordinario.
\beginsong{Santo il Signore}[by={D. Ricci}]
\beginverse*
Santo il Signore Dio dell'universo,
i cieli e la terra son pieni della tua gloria.
Osanna nell'alto dei cieli e benedetto colui che viene
nel nome del Signore. Osanna nell'alto dei cieli.
\endverse
\beginverse*
È santo, santo, santo il Signore della vita.
È santo, santo, santo il Signore della storia,
nel cielo e sulla terra risplende la sua gloria.
Osanna, osanna, osanna nel cielo sconfinato.
Osanna, osanna, osanna nel canto del creato.
E sia benedetto colui che viene in nome di Dio. \rep{3}
\endverse
\endsong

% Fonte: LibrettoAccordi.pdf (Canti Coro), H2.
% Nell'elenco: Santo, Ordinario.
\beginsong{Santo, santo, santo}[by={D. Ricci}]
\beginverse*
Santo! Santo! Santo
è il Signore Dio dell'universo!
Santo! Santo! Santo!
I cieli e la terra sono pieni della tua gloria!
Osanna! Osanna! Osanna nell'alto dei cieli!
Benedetto colui che viene nel nome del Signore.
Osanna! Osanna! Osanna nell'alto dei cieli! \rep{3}
\endverse
\endsong

% Fonte: LibrettoAccordi.pdf (Canti Coro), H7.
% Nell'elenco: Santo, Ordinario. CD: Xaipe.
\beginsong{Santo (Xaipe)}[by={M. Balduzzi (CD Xaipe)}]
\beginchorus
Osanna! Osanna in excelsis! Osanna!
\endchorus
\beginverse*
Sanctus Dominus Deus Sabaoth! \rep{2}
Pleni, pleni sunt coeli et terra \rep{2}
gloria tua.
\endverse
\beginverse*
Benedictus qui venit \rep{2}
in nomine Domini.
\endverse
\endsong

\end{songs}

% ===================== ALLELUIA E LODE =====================
% !TEX root = ../libretto.tex
% Sezione: Alleluia e Lode (acclamazioni al Vangelo, anche quaresimali)
\momento{Alleluia e Lode}
\begin{songs}{titoli}

% Fonte: LibrettoAccordi.pdf (Canti Coro), E5.
% Nell'elenco: Alleluia, Matrimonio.
\beginsong{Alleluia, beati coloro}[by={F. Baggio, F. Buttazzo}]
\beginchorus
Alleluja! Alleluja! Alleluja! Alleluja! Alleluja! Alleluja! \rep{2}
\endchorus
\beginverse
Beati coloro che sanno ascoltare quello che dice lo Sposo alla Chiesa.
Cristo Signore rivela il suo amore per la sua sposa, la Chiesa.
\endverse
\endsong

% Fonte: LibrettoAccordi.pdf (Canti Coro), E4.
% Nell'elenco: Alleluia, Ordinario.
\beginsong{Alleluia, canto per Cristo}[by={E. Costa}]
\beginchorus
Alleluja! Alleluja! Alleluja! Alleluja! Alleluja! \rep{2}
\endchorus
\beginverse
Canto per Cristo che mi libererà, quando verrà nella gloria,
quando la vita con lui rinascerà. Alleluja! Alleluja!
\endverse
\beginverse
Canto per Cristo: in Lui rifiorirà ogni speranza perduta,
ogni creatura con lui risorgerà. Alleluja! Alleluja!
\endverse
\beginverse
Canto per Cristo: un giorno tornerà! Festa per tutti gli amici,
festa di un mondo che più non morirà. Alleluja! Alleluja!
\endverse
\endsong

% Fonte: LibrettoAccordi.pdf (Canti Coro), E2.
% Nell'elenco: Alleluia, Pentecoste; l'utente lo vuole tra gli Alleluia.
\beginsong{Alleluia (Come fuoco vivo)}[by={Gen Rosso, Gen Verde (CD Come fuoco vivo)}]
\beginverse*
Alleluja! Alleluja! Alleluja!
Alleluja! Alleluja! Alleluja! Alleluja!
Ah, alleluja! Ah, alleluja, alleluja!
Alleluja! Alleluja! Alleluja!
Alleluja! Alleluja! Alleluja! Alleluja! Alleluja!
\endverse
\endsong

% Fonte: LibrettoAccordi.pdf (Canti Coro), E12.
% Nell'elenco: Alleluia, Ordinario.
\beginsong{Alleluia\soloindice{ (Gen Verde)}}[by={Gen Verde}]
\beginverse*
Alleluja! Alleluja! Alleluja! Alleluja! \rep{2}
Alleluja! Alleluja! Alleluja! Ah!
Alleluja! Alleluja! Alleluja! Alleluja! Alleluja! Alleluja!
Alleluja! Alleluja! Alleluja! Alleluja! \rep{2}
\endverse
\endsong

% Fonte: LibrettoAccordi.pdf (Canti Coro), E11. Omesse le indicazioni dei due gruppi.
% Nell'elenco: Alleluia, Ordinario.
\beginsong{Alleluia, lode cosmica}[by={S. Puri}]
\beginchorus
Alleluja! Alleluja! Alleluja! Alleluja!
Alleluja! Alleluja! Alleluja! Alleluja! \rep{2}
\endchorus
\beginverse
Lodino il Signor i cieli, lodino il Signor i mari,
gli angeli, i cieli dei cieli: il suo nome è grande e sublime.
Lode, lode al suo nome.
Sole, luna e stelle ardenti, neve, pioggia, nebbia e fuoco
lodino il suo nome in eterno!
Sia lode al Signor! Sia lode al Signor!
Sia lode al Signore!
\endverse
\beginverse
Lodino il Signor le terre, lodino il Signor i monti,
il vento della tempesta che obbedisce alla sua voce.
Giudici, sovrani tutti, giovani, fanciulle, vecchi
lodino il suo nome in eterno!
Sia lode al Signor! Sia lode al Signor!
Sia lode al Signore!
\endverse
\endsong

% Fonte: LibrettoAccordi.pdf (Canti Coro), E3.
% Nell'elenco: Alleluia, Ordinario.
\beginsong{Alleluia\soloindice{ (OMI)}}[by={OMI}]
\beginverse*
Alleluja! Alleluja!
Alleluja! Alleluja! Alleluja! \rep{2}
\endverse
\endsong

% Fonte: LibrettoAccordi.pdf (Canti Coro), E6.
% Nell'elenco: Alleluia, Ordinario.
\beginsong{Alleluia, ta parole}[by={F. Baggio}]
\beginchorus
Alle-alleluja! Alleluja. Alle-alleluja! Alleluja. \rep{2}
\endchorus
\beginverse
È la tua parola, la verità, che dona luce a questa umanità:
rivela all'uomo il cuore di Dio e porta al mondo unità.
\endverse
\endsong

% Fonte: LibrettoAccordi.pdf (Canti Coro), E9.
% Nell'elenco: Alleluia, Ordinario.
\beginsong{Alleluia\soloindice{ (Verbum Panis)}}[by={M. Balduzzi (CD Verbum Panis)}]
\beginchorus
Alleluja! Alleluja! Alleluja! Alleluja! Alleluja! \rep{2}
\endchorus
\beginverse*
Alleluja! Alleluja! \rep{4}
\endverse
\endsong

% Fonte: LibrettoAccordi.pdf (Canti Coro), E7.
% Nell'elenco: Alleluia (acclamazione), Quaresima. Spostato in "Alleluia e Lode" su indicazione dell'utente.
\beginsong{Lode a te, o Cristo\soloindice{ (F. Buttazzo)}}[by={F. Buttazzo}]
\beginchorus
Lode a Te, o Cristo, Verbo di salvezza!
Lode a Te, o Cristo, Parola di Dio per noi! \rep{2}
\endchorus
\beginverse
Come la luce illumina ogni cosa,
la tua Parola rivela a noi il tuo volto.
Come la pioggia scende nella terra,
la tua Parola entra nel nostro cuore.
\endverse
\beginverse
Seme da te piantato in terra buona,
la tua Parola in noi porterà i suoi frutti.
Pane che tu ci doni in ogni giorno,
la tua Parola è cibo per noi di vita.
\endverse
\endsong

% Fonte: LibrettoAccordi.pdf (Canti Coro), E13.
% Nell'elenco: Alleluia (acclamazione), Quaresima. Spostato in "Alleluia e Lode" su indicazione dell'utente.
\beginsong{Lode a te, o Cristo\soloindice{ (Gen Verde)}}[by={Gen Verde}]
\beginchorus
Lode a Te, o Cristo,
Re di eterna gloria. \rep{2}
\endchorus
\beginverse
Signore Tu sei veramente
il Salvatore del mondo.
Dammi dell'acqua viva
perché non abbia più sete.
\endverse
\beginverse
Chi beve di quest'acqua
avrà di nuovo sete.
Ma chi beve dell'acqua che io gli darò
non avrà mai più sete.
\endverse
\endsong

% Fonte: LibrettoAccordi.pdf (Canti Coro), E10.
% Nell'elenco: Alleluia (acclamazione), Quaresima. Spostato in "Alleluia e Lode" su indicazione dell'utente.
\beginsong{Lode e gloria}[by={P. Spoladore}]
\beginchorus
Lode e gloria a Te, o Signore! \rep{2}
\endchorus
\beginverse
Date lode al Signore o figli di Dio,
benedite il suo nome su tutta la terra.
Il suo amore per noi è fedele per sempre,
il suo amore per noi è per sempre.
\endverse
\beginverse
Acclamate al Signore da tutta la terra,
ci ha creati e voluti, a lui apparteniamo.
Il Signore ci guida e ci tiene per mano,
senza fine è il suo amore per noi.
\endverse
\endsong

\end{songs}
